\chapter{2017:Jacques Dubochet and Joachim Frank and Richard Henderson}

\label{chap:chemistry2017}

The Nobel Prize in Chemistry for the year 2017 was awarded jointly to Jacques Dubochet and Joachim Frank and Richard Henderson.

\textbf{Motivation:}

for developing cryo-electron microscopy for the high-resolution structure determination of biomolecules in solution

\section{Jacques Dubochet}

\subsection*{Early Life and Education}

Jacques Dubochet was born in 1942 in Aigle, Switzerland.

\subsection*{Career and Major Research}

Dubochet studied at the University of Lausanne and received a doctorate in biophysics from the University of Geneva and the University of Basel in 1973. After work at the European Molecular Biology Laboratory in Heidelberg he became a professor at the University of Lausanne. He developed the method of vitrifying water by cooling it so fast that it solidifies without forming ice crystals, preserving biological molecules in their natural state in the electron microscope.

\subsection*{Nobel Prize-winning Work}

The work for which Jacques Dubochet was awarded the Nobel Prize in Chemistry in 2017 was described by the Nobel Committee as: "for developing cryo-electron microscopy for the high-resolution structure determination of biomolecules in solution".

Dubochet's vitrified water made it possible to image intact biomolecules in a thin layer of glass-like ice under the electron beam.

\subsection*{Significance and Impact}

His plunge-freezing method protects delicate biological specimens from damage and dehydration, and it is an essential step in every modern cryo-EM experiment.

\subsection*{Later Career and Honors}

Dubochet is professor emeritus at the University of Lausanne and has spoken widely about science and society.

\subsection*{Legacy}

Water vitrification is the technical foundation of cryo-electron microscopy, now a routine method for determining biomolecular structures.

\subsection*{Selected Publications}

"Cryo-Electron Microscopy of Viruses" (Nature, 1984)

\clearpage

\section{Joachim Frank}

\subsection*{Early Life and Education}

Joachim Frank was born in 1940 in Siegen, Germany.

\subsection*{Career and Major Research}

Frank studied physics and received his doctorate from the Technical University of Munich in 1970. He became a professor at Columbia University, where he developed the computational methods that combine noisy electron microscope images of single molecules into three-dimensional structures.

\subsection*{Nobel Prize-winning Work}

The work for which Joachim Frank was awarded the Nobel Prize in Chemistry in 2017 was described by the Nobel Committee as: "for developing cryo-electron microscopy for the high-resolution structure determination of biomolecules in solution".

Frank's image-processing procedures aligned and averaged thousands of individual molecular images, reconstructing detailed three-dimensional structures of machines such as the ribosome.

\subsection*{Significance and Impact}

His software made it possible to extract high-resolution information from low-dose images of molecules that cannot be crystallized.

\subsection*{Later Career and Honors}

Frank continues his work on image processing and the structure of the ribosome at Columbia University.

\subsection*{Legacy}

Frank's computational approach is the core of single-particle cryo-EM, which has revolutionized structural biology.

\subsection*{Selected Publications}

"SPIDER and WEB: Processing and Visualization of Images in 3D Electron Microscopy and Related Fields" (Journal of Structural Biology, 1996)

\clearpage

\section{Richard Henderson}

\subsection*{Early Life and Education}

Richard Henderson was born in 1945 in Edinburgh, Scotland.

\subsection*{Career and Major Research}

Henderson earned a bachelor's degree from the University of Edinburgh and a doctorate from the University of Cambridge in 1969. He became a group leader at the MRC Laboratory of Molecular Biology in Cambridge, where in 1975 he and Nigel Unwin produced the first three-dimensional structure of a membrane protein, bacteriorhodopsin, using electron microscopy.

\subsection*{Nobel Prize-winning Work}

The work for which Richard Henderson was awarded the Nobel Prize in Chemistry in 2017 was described by the Nobel Committee as: "for developing cryo-electron microscopy for the high-resolution structure determination of biomolecules in solution".

Henderson showed that electron microscopy could reach atomic resolution by improving the detectors and sample preparation of cryo-EM, and he demonstrated this by imaging a protein at near-atomic resolution.

\subsection*{Significance and Impact}

His work proved that cryo-EM could determine structures as accurately as X-ray crystallography, a milestone that established the method for routine use.

\subsection*{Later Career and Honors}

Henderson was a long-serving group leader at the MRC Laboratory of Molecular Biology and helped lead the development of direct electron detectors.

\subsection*{Legacy}

Henderson's demonstration that cryo-EM reaches atomic resolution transformed structural biology and earned the method widespread adoption.

\subsection*{Selected Publications}

"Three-Dimensional Model of Purple Membrane Obtained by Electron Microscopy" (Nature, 1975)

\clearpage

\section*{Chapter Summary}

The 2017 prize recognized the development of cryo-electron microscopy, which determines the three-dimensional structure of biomolecules in solution without crystallization. Dubochet contributed the vitrification of water, Frank the image processing, and Henderson the path to atomic resolution, together making cryo-EM a cornerstone of structural biology.
